\uuid{Bf5x}
\exo7id{6305}
\titre{exo7 6305}
\auteur{mayer}
\organisation{exo7}
\datecreate{2011-10-16}
\isIndication{false}
\isCorrection{false}
\chapitre{Différentielle d'ordre supérieur, formule de Taylor}
\sousChapitre{Différentielle d'ordre supérieur, formule de Taylor}
\cpge{mpsi-mp2i.fonctions-deux-variables.b, pcsi.fonctions-deux-variables.b, ptsi.fonctions-deux-variables.b, mpsi-mp2i.fonctions-deux-variables.c, pcsi.fonctions-deux-variables.c, ptsi.fonctions-deux-variables.c, bcpst1.fonctions-deux-variables.c, bcpst1.fonctions-deux-variables.d, mp-mpi.calcul-differentiel-optimisation.a, mp-mpi.calcul-differentiel-optimisation.b, pc.fonctions-plusieurs-variables.c, psi.fonctions-plusieurs-variables.c}
\module{Calcul différentiel}
\niveau{L3}
\difficulte{}

\contenu{
\question{
Soit $f(x,y)$ une fonction de classe $C^2$ au
voisinage du cercle $x^2 + y^2 = 1$. On pose ${\partial
f\over \partial x} (1,0) = a$ et ${\partial^2 f\over
\partial y^2} (1,0) = b$. Pour tout nombre réel $\theta
$, soit $F(\theta ) = f(\cos \theta, \sin \theta )$.
Calculez $F''(0)$ en fonction de $a$ et $b$.
}
}
